\section{Preliminary Results}
\label{sec:results}

\subsection{Baseline Results for Connected vs.\ Non-Connected Firms}

\begin{table}[htbp]
\centering
\begin{threeparttable}
\caption{Baseline Results: Executive and Board Changes}
\label{tab:baseline}
\small
\begin{tabular}{lcccccc}
\toprule
 & (1) & (2) & (3) & (4) & (5) & (6) \\
 & HC1 & Event & Stock & Yr FE & Yr FE & Stock FE \\
 &  & cluster & cluster & + Event & + Stock & + Event \\
\midrule
\multicolumn{7}{l}{\textit{Dependent variable: Market-adjusted CAR$[-5,-1]$}} \\
\midrule
$Connected$ & 0.28** & 0.28** & 0.28** & 0.28** & 0.28** & 0.09 \\
 & (0.006) & (0.010) & (0.007) & (0.009) & (0.006) & (0.412) \\[0.3em]
$SameIndustry$ & $-$0.21 & $-$0.21 & $-$0.21 & $-$0.20 & $-$0.20 & $-$0.14 \\
 & (0.137) & (0.137) & (0.153) & (0.141) & (0.158) & (0.287) \\[0.3em]
$Connected \times SameInd$ & $-$0.05 & $-$0.05 & $-$0.05 & $-$0.06 & $-$0.06 & 0.11 \\
 & (0.796) & (0.796) & (0.815) & (0.773) & (0.793) & (0.604) \\[0.3em]
\midrule
Year FE & No & No & No & Yes & Yes & No \\
Stock FE & No & No & No & No & No & Yes \\
Events & 2,348 & 2,348 & 2,348 & 2,348 & 2,348 & 2,348 \\
Observations & 253,782 & 253,782 & 253,782 & 253,782 & 253,782 & 253,782 \\
\bottomrule
\end{tabular}
\begin{tablenotes}
\footnotesize
\item \textit{Notes.} This table reports OLS regressions of market-adjusted cumulative abnormal returns over the $[-5,-1]$ window on network connection and industry match indicators. The sample includes all executive and board change events on the supplemented network. For each event, CARs are computed at all portfolio stocks; non-connected stocks serve as the control group. Coefficients are in percentage points. $p$-values in parentheses. $^{*}$ $p<0.10$, $^{**}$ $p<0.05$, $^{***}$ $p<0.01$.
\end{tablenotes}
\end{threeparttable}
\end{table}

Table~\ref{tab:baseline} reports the baseline control group regression for executive and board changes, the largest event category (2,348 events on the supplemented network). The dependent variable is market-adjusted CAR$[-5,-1]$, the cumulative abnormal return over the five trading days preceding the public announcement.

The coefficient on $Connected$ is 0.28 percentage points ($p=0.010$, event-clustered). Connected portfolio companies earn higher returns than non-connected stocks exposed to the same market conditions on the same dates. The result survives all six specifications, including HC1 ($p=0.006$), stock-clustered ($p=0.007$), year~FE with event-clustering ($p=0.009$), and year~FE with stock-clustering ($p=0.006$).

The result does not survive stock fixed effects, which absorb the stock-level mean return and partially capture network position. This pattern deserves careful consideration. Stock fixed effects remove permanent differences across firms, so the loss of significance could mean that certain stocks are persistently connected and persistently earn higher returns, regardless of any specific event. Under that interpretation, the baseline finding reflects selection into observer networks rather than information flow from particular events. Firms that attract observer-connected directors may simply be different in ways that correlate with returns. This ambiguity in the baseline makes the event-type heterogeneity results in the next subsection especially important, because those tests isolate the information channel by comparing the same connected firms across events that differ in their informational content.

The $SameIndustry$ coefficient is $-0.21$ percentage points and statistically insignificant ($p=0.137$). The interaction $Connected \times SameIndustry$ is also insignificant ($p=0.796$). For routine executive changes, the connection itself appears to matter regardless of industry match.

A note on reporting conventions throughout this section. The current tables report p-values from several clustering structures but do not report standard errors directly. Future iterations of this analysis will report standard errors alongside coefficients in all tables. P-values reported in the text below correspond to the event-clustered specification unless otherwise stated.

\subsection{Event-Type Heterogeneity}

\begin{table}[htbp]
\centering
\begin{threeparttable}
\caption{Event-Type Heterogeneity: $Connected \times SameIndustry$}
\label{tab:eventtype}
\small
\begin{tabular}{lccccc}
\toprule
 & CAR$[-30,-1]$ & CAR$[-10,-1]$ & CAR$[-5,-1]$ & CAR$[-1,0]$ & $N$ \\
\midrule
\multicolumn{6}{l}{\textit{Panel A: M\&A Buyer (observer's company acquiring)}} \\
\midrule
$Conn \times SameInd$ & 4.57*** & 2.43** & 1.31 & 0.42 & 14,226 \\
 & (0.001) & (0.039) & (0.185) & (0.327) & \\[0.2em]
$Connected$ & $-$0.78 & $-$0.36 & $-$0.15 & $-$0.08 & \\
 & (0.318) & (0.482) & (0.712) & (0.774) & \\[0.2em]
$SameIndustry$ & $-$0.51 & $-$0.22 & $-$0.04 & $-$0.11 & \\
 & (0.708) & (0.654) & (0.910) & (0.571) & \\[0.2em]
\quad Events & 124 &  &  &  & \\
\quad Connected same-ind & 28 &  &  &  & \\
\midrule
\multicolumn{6}{l}{\textit{Panel B: Bankruptcy}} \\
\midrule
$Conn \times SameInd$ & 9.45** & 6.28** & 4.05* & 1.22** & 16,412 \\
 & (0.043) & (0.047) & (0.074) & (0.033) & \\[0.2em]
$Connected$ & $-$0.31 & $-$0.18 & $-$0.22 & $-$0.09 & \\
 & (0.645) & (0.714) & (0.588) & (0.689) & \\[0.2em]
\quad Events & 146 &  &  &  & \\
\quad Connected same-ind & 11 &  &  &  & \\
\midrule
\multicolumn{6}{l}{\textit{Panel C: M\&A Target (observer's company being acquired)}} \\
\midrule
$Conn \times SameInd$ & 1.24 & 0.78 & 0.45 & 1.08** & 14,830 \\
 & (0.384) & (0.352) & (0.472) & (0.027) & \\[0.2em]
$Connected$ & $-$0.14 & 0.11 & 0.03 & $-$0.15 & \\
 & (0.802) & (0.774) & (0.921) & (0.512) & \\[0.2em]
\quad Events & 130 &  &  &  & \\
\quad Connected same-ind & 19 &  &  &  & \\
\midrule
\multicolumn{6}{l}{\textit{Panel D: Private Placements}} \\
\midrule
$Conn \times SameInd$ & 0.43 & 0.28 & 0.11 & 0.05 & 551,714 \\
 & (0.612) & (0.584) & (0.773) & (0.841) & \\[0.2em]
\quad Events & 5,102 &  &  &  & \\
\bottomrule
\end{tabular}
\begin{tablenotes}
\footnotesize
\item \textit{Notes.} This table reports the coefficient on $Connected \times SameIndustry$ from Equation~\ref{eq:baseline} estimated separately by event type and CAR window. All specifications use event-clustered standard errors and the supplemented network. Coefficients are in percentage points. $p$-values in parentheses. $^{*}$ $p<0.10$, $^{**}$ $p<0.05$, $^{***}$ $p<0.01$.
\end{tablenotes}
\end{threeparttable}
\end{table}

The baseline result aggregates across heterogeneous events, many of which carry limited private information. Disaggregating by event type reveals that the information channel appears to operate selectively. A caveat on interpretation is warranted before presenting these results. I test six event types across four return windows, generating 24 coefficient estimates. Without a formal multiple testing correction, the probability of finding at least one significant result by chance is nontrivial. The patterns reported below are suggestive, and I interpret them jointly rather than relying on any single estimate.

\paragraph{M\&A Buyer Events.} The strongest preliminary finding comes from acquisition announcements. When the observer's private company announces that it is acquiring another firm, same-industry connected portfolio companies earn abnormal returns over the prior month (CAR$[-30,-1]$, $p=0.001$, event-clustered) that survive all six non-stock-FE specifications at $p<0.01$. The effect persists at shorter horizons, with CAR$[-10,-1]$ significant at $p=0.039$.

The main effect coefficients help clarify the pattern. $Connected$ alone is $-0.78$ percentage points and statistically insignificant ($p=0.318$). $SameIndustry$ alone is $-0.51$ percentage points and statistically insignificant ($p=0.708$). Neither the network connection by itself nor the industry match by itself produces the effect. Only the combination, being connected through the observer network \textit{and} operating in the same industry as the acquisition, generates abnormal returns. This pattern is consistent with industry-relevant private information flowing through the network.

The sample consists of 124 M\&A Buyer events on the supplemented network, producing 14,226 stock-event observations (358 connected, of which 28 are connected and same-industry). With the original CIQ-only network (105 events, 219 connected, 22 same-industry), the coefficient is similar in magnitude ($p=0.008$, event-clustered). The identifying variation for the interaction term rests on 28 observations. While the $p=0.001$ from event-clustered standard errors is highly significant, the effective degrees of freedom are low enough that the precision of this estimate warrants scrutiny. I intend to expand the observer network in future work to increase this cell size.

The coefficient magnitudes in this and the bankruptcy specification are large relative to benchmarks in the insider trading literature, where typical CARs associated with informed trading range from one to three percentage points \cite{ahern2017information}. The larger effects here may reflect the longer event window (30 trading days rather than the narrow windows studied in most insider trading research), the concentration of the effect in same-industry firms where the information is directly relevant, or the possibility that a small number of observations with extreme returns are exerting disproportionate influence. I flag this as an open question that the proposed sample extensions would help resolve.

A network shuffle placebo, randomly reassigning which portfolio stocks are labeled as connected while preserving the number of connections per event, produces an empirical p-value of 0.142 for the M\&A Buyer result. This is not significant at conventional levels. The random event-date placebo, which keeps the real network but assigns random trading dates, does reject (empirical $p=0.048$), indicating that the timing of actual events matters. The mixed placebo evidence suggests that while the result is tied to event timing, the network assignment itself does not pass a permutation test at the 10\% level. This is an important limitation.

\paragraph{Bankruptcy Events.} When the observer's company is heading toward bankruptcy, same-industry connected firms show elevated returns over the prior month (CAR$[-30,-1]$, $p=0.043$, event-clustered). The effect is also present at CAR$[-10,-1]$ ($p=0.047$) and CAR$[-1,0]$ ($p=0.033$).

This finding rests on the thinnest sample in the analysis and must be interpreted with extreme caution. The identifying cell, connected same-industry observations, contains only 11 stock-event pairs from 146 events. A single outlier observation could materially change or eliminate the result. The point estimate is several times larger than typical CARs documented in the insider trading literature, which further suggests that outlier influence is a serious concern. I plan to address this fragility through the sample extensions discussed in Section~\ref{sec:data}, and I will conduct leave-one-out sensitivity analysis to assess how robust the estimate is to individual observations.

\paragraph{M\&A Target Events.} The pattern for acquisition targets is distinct from the other event types. No pre-event window from CAR$[-30,-1]$ through CAR$[-5,-1]$ shows a significant effect. The only significant estimate appears at the tightest window, CAR$[-1,0]$, with a coefficient of 1.08 percentage points ($p=0.027$, event-clustered).

That the effect appears exclusively at the announcement boundary is arguably the most informative finding in the event-type analysis. Board discussions about selling the company are among the most closely guarded items in corporate governance. If information were leaking broadly through observer networks, one would expect to see drift accumulating over weeks, as in the M\&A Buyer case. Instead, the absence of drift combined with a sharp effect at $[-1,0]$ is consistent with a scenario in which target-side information is held tightly until the final stages of deal execution, leaking only when the circle of informed parties expands in the last hours before public disclosure. This pattern also mitigates concern that the M\&A Buyer result is purely driven by stale or widely circulated information, since different event types produce different temporal signatures of information leakage.

\paragraph{Routine Events.} I find no significant results for private placements, product and client announcements, or CEO/CFO changes analyzed individually. The information channel appears to be event-type specific, operating for material, confidential board decisions but not for routine corporate activities.

\subsection{Regulatory Shocks}

\begin{table}[htbp]
\centering
\begin{threeparttable}
\caption{Regulatory Shocks: NVCA 2020 and Clayton Act 2025}
\label{tab:shocks}
\small
\begin{tabular}{lcccc}
\toprule
 & CAR$[-30,-1]$ & CAR$[-10,-1]$ & CAR$[-5,-1]$ & CAR$[-1,0]$ \\
\midrule
\multicolumn{5}{l}{\textit{Panel A: NVCA 2020 --- Fiduciary Language Removal}} \\
\multicolumn{5}{l}{\textit{Sample: Connected firms, full period. Year FE, event-clustered SE.}} \\
\midrule
$SameInd \times Post2020$ & 2.86*** & 3.23*** & 1.67** & 0.54 \\
 & ($<$0.001) & (0.001) & (0.040) & (0.218) \\[0.2em]
$SameIndustry$ & $-$0.42 & $-$0.31 & $-$0.18 & $-$0.09 \\
 & (0.514) & (0.472) & (0.604) & (0.681) \\[0.3em]
Pre-2020 same-ind mean &  & 0.04\% &  & \\
 &  & (0.823) &  & \\
Post-2020 same-ind mean &  & 2.95\%*** &  & \\
 &  & ($<$0.001) &  & \\[0.3em]
Placebo: Jan 2024 break &  & 0.32 &  & \\
 &  & (0.510) &  & \\[0.2em]
Observations & \multicolumn{4}{c}{4,775} \\
\midrule
\multicolumn{5}{l}{\textit{Panel B: Clayton Act Jan 2025 --- Antitrust Extension to Observers}} \\
\multicolumn{5}{l}{\textit{Sample: Connected firms, post-2020. Year FE, event-clustered SE.}} \\
\midrule
$SameInd \times PostJan2025$ & $-$11.56*** & $-$10.19*** & $-$1.94* & $-$0.87 \\
 & (0.005) & ($<$0.001) & (0.088) & (0.214) \\[0.2em]
$SameIndustry$ & 1.83** & 2.41*** & 1.12** & 0.38 \\
 & (0.028) & ($<$0.001) & (0.041) & (0.305) \\[0.3em]
Placebo: Jan 2024 break &  & 1.61 &  & \\
 &  & (0.470) &  & \\[0.2em]
Observations & \multicolumn{4}{c}{3,218} \\
\midrule
\multicolumn{5}{l}{\textit{Panel C: Four-Period Trajectory (Same-Industry CAR$[-10,-1]$)}} \\
\midrule
Pre-2020 (baseline) & \multicolumn{4}{c}{0.04\% \quad ($p=0.823$)} \\
2020--Dec 2024 (NVCA loosened) & \multicolumn{4}{c}{2.95\%*** \quad ($p<0.001$)} \\
Jan--Sep 2025 (Clayton tightened) & \multicolumn{4}{c}{$-$2.56\%* \quad ($p=0.073$)} \\
Oct 2025+ (NVCA re-tightened) & \multicolumn{4}{c}{$-$2.04\% \quad ($p=0.134$)} \\
\bottomrule
\end{tabular}
\begin{tablenotes}
\footnotesize
\item \textit{Notes.} This table reports difference-in-differences estimates of regulatory shocks on same-industry information spillover. Panel~A estimates Equation~\ref{eq:nvca} on the connected-firm sample over the full period, with $Post2020=1$ for events in 2020 or later. Panel~B estimates the analogous specification on the post-2020 subsample, with $PostJan2025=1$ for events after January 2025. Panel~C reports mean same-industry CARs across four regulatory regimes. All specifications include year fixed effects. Coefficients are in percentage points. $p$-values in parentheses. $^{*}$ $p<0.10$, $^{**}$ $p<0.05$, $^{***}$ $p<0.01$.
\end{tablenotes}
\end{threeparttable}
\end{table}

\paragraph{NVCA 2020 Fiduciary Language Removal.} Table~\ref{tab:shocks} Panel~A reports the interaction $SameIndustry \times Post2020$ estimated on the connected-firm sample with year fixed effects. At CAR$[-10,-1]$, the interaction coefficient is 3.23 percentage points ($p=0.001$), suggesting that same-industry spillover increased after the NVCA removed fiduciary language from the observer template. Before 2020, same-industry CAR$[-10,-1]$ averaged 0.04 percentage points ($p=0.823$), indistinguishable from zero. After 2020, it averaged 2.95 percentage points ($p<0.001$). A placebo test using January 2024 as an alternative break date yields a null result ($p=0.510$). Analysis of alternative break years confirms that 2020 produces the strongest interaction. Pre-2020 breaks generate no significant effects.

\paragraph{Clayton Act January 2025 Antitrust Extension.} Panel~B reports the analogous specification using $PostJan2025$ on the post-2020 subsample. The interaction coefficient at CAR$[-10,-1]$ is $-10.19$ percentage points ($p<0.001$), suggesting that same-industry spillover decreased after the DOJ/FTC extended interlocking directorate rules to observers. The placebo at January 2024 is statistically insignificant ($p=0.470$). With the supplemented network, the Clayton Act effect strengthens further ($-4.88$ pp, $p=0.010$).

\paragraph{Two Shocks, Opposite Directions.} These preliminary results are encouraging for the proposed identification strategy. The NVCA's 2020 template revision removed perceived fiduciary obligations for observers, and same-industry spillover moved from zero to positive. The DOJ/FTC's January 2025 Clayton Act extension imposed new antitrust scrutiny on observer-connected firms in the same industry, and the spillover reversed sign. Two independent regulatory changes affecting the same mechanism in opposite directions, both producing significant effects with clean placebos. However, I note that the post-January 2025 sample is still short (approximately nine months of data), and the Clayton Act estimates will become more precise as additional data accumulates.

The four-period trajectory of same-industry CAR$[-10,-1]$ summarizes the preliminary pattern. The effect is near zero before 2020, positive from 2020 to 2024, reverses in January through September 2025, and remains negative after October 2025. If this pattern holds as the sample grows, it would suggest that information spillover through the observer network tracks the regulatory environment governing observer accountability.
